% Vidruma -- mid-submission report.  Compile:  xelatex main && bibtex main && xelatex main && xelatex main
% Tables in tables/ are generated by make_midsem_tables.py from results/.
\documentclass[11pt]{article}
\usepackage[final]{acl}
\usepackage{fontspec}
\setmainfont{TeX Gyre Termes}
\usepackage{amsmath,amssymb}
\usepackage{booktabs,array}
\usepackage{microtype}
\usepackage{enumitem}
\usepackage[section]{placeins}
\setlist{nosep,leftmargin=*}

\newcommand{\vk}[1]{\textsf{#1}}
\newcommand{\OOV}{\operatorname{OOV}}
\newcommand{\PMI}{\operatorname{PMI}}
\newcommand{\code}[1]{\texttt{\small #1}}
\newcommand{\ex}{\par\smallskip\noindent\textsc{Running example.}~}

\title{Vidruma: From Varga to Vector Space\\[2pt]
\large Mid-Submission Report}

\author{Sai Sreemayee Sistla \\ 2025114013 \And Vinayak Sharma \\ 2025114002}

\begin{document}
\maketitle

%==========================================================================
\begin{abstract}
We study distributional semantics for Sanskrit at the level of lemmas and
lemma IDs: four embedding models (word2vec, FastText, GloVe, PPMI-SVD)
trained on the Digital Corpus of Sanskrit (DCS) and evaluated on the EvalSan
benchmarks and the Amarakośa. This report covers the data and evaluation
work and the first full results. An audit of the setting found that
EvalSan's relatedness and synonym files are in a different transliteration
(WX) from its vectors (SLP1), that the benchmarks contain self-pairs,
duplicates and doubly-correct items, and that DCS contains the Amarakośa
itself. On a leak-free corpus and a cleaned benchmark, with every design
decision taken on dev data, lemmatisation improves every model, lemma IDs do
not, and PPMI-SVD is the strongest model. Under EvalSan's own protocol our
models exceed the published static baselines on three of four tasks.
Coverage, not technique, is the main limit; machine annotation of 10M more
words of Sanskrit is under way.
\end{abstract}

%==========================================================================
\section{Introduction}
%==========================================================================
Word-embedding models for Sanskrit report out-of-vocabulary (OOV) rates of up
to 55\% on their own benchmarks \citep{sandhan2023evalsan}. Sanskrit inflects
heavily, fuses words at their boundaries (\emph{sandhi}) and compounds
freely, so a model faces far more distinct strings than meanings. We ask:
\begin{description}
\item[RQ1] How does lemmatisation, and \emph{ID-based} lemmatisation that
  separates homographs, affect embeddings on Sanskrit benchmarks?
\item[RQ2] Do clusters in the embedding space recover the Amarakośa's
  \emph{vargas}, and what do the divergences mean?
\item[RQ3] How much of the OOV rate is inflection rather than sparsity, and
  how much ambiguity does lemma-level disambiguation resolve?
\end{description}
The grammatical tradition frames RQ3: Bhartṛhari
\citep{iyer1965vakyapadiya} names association (\emph{saṃsarga}),
co-presence (\emph{sāhacarya}) and topic (\emph{prakaraṇa}) among the
factors that select a word's meaning, which is what a co-occurrence window
observes \citep{harris1954,firth1957,matilal1990word}.

\paragraph{Status.} We are at about half of the project. Done: the data
pipeline and normalisation, the benchmark audit and cleaned benchmark, the
leak analysis, four models, the evaluation protocol, a dev-set decision log
and the full test grid (RQ1). RQ2 and RQ3 have preliminary results
(\S\ref{sec:prelim}); \S\ref{sec:next} lists the remaining work.

%==========================================================================
\section{Setup, with a running example}\label{sec:setup}
%==========================================================================
Each new notion is illustrated on the toy corpus of Table~\ref{tab:toy};
\code{docs/toy\_example.py} reproduces every toy number. Logarithms are base
2.

\begin{table}[t]\centering\footnotesize
\setlength{\tabcolsep}{4pt}
\begin{tabular}{rll}
\toprule
\# & as written (\vk{surf}) & lemmas (\vk{lem})\\
\midrule
1 & \emph{jalaṃ pibati} & \emph{jala pā}\\
2 & \emph{toyaṃ pibati} & \emph{toya pā}\\
3 & \emph{nadyāṃ jalam} & \emph{nadī jala}\\
4 & \emph{nadyāṃ toyam} & \emph{nadī toya}\\
5 & \emph{agnirdahati} & \emph{agni dah}\\
6 & \emph{vahnirdahati} & \emph{vahni dah}\\
7 & \emph{dīpe'gniḥ} & \emph{dīpa agni}\\
8 & \emph{dīpe vahniḥ} & \emph{dīpa vahni}\\
9 & \emph{jalamagniśca} & \emph{jala agni ca}\\
10 & \emph{toyaṃ vahniśca} & \emph{toya vahni ca}\\
\bottomrule
\end{tabular}
\caption{Running example. Gold: synsets \{\emph{jala, toya, salila}\}
``water'', \{\emph{agni, vahni}\} ``fire''; vargas \{\emph{jala, toya,
salila, nadī}\} and \{\emph{agni, vahni, dīpa}\}.}
\label{tab:toy}
\end{table}

\paragraph{Views.} A \emph{view} maps every word occurrence in DCS
(5.80M words, 770,945 sentences) to a token. We use \vk{surf} (written,
sandhied), \vk{seg} (sandhi resolved), \vk{lem} (lemma strings), \vk{lid}
(lemma IDs, which separate homographs) and \vk{sense} (sense IDs, where
annotated); each undoes one source of variation.

\ex Sentence 7 is written \emph{dīpe'gniḥ}: sandhi has absorbed the initial
\emph{a} of \emph{agniḥ}, so \vk{surf} never sees \emph{agni}. \vk{seg}
restores \emph{dīpe agniḥ} and \vk{lem} gives \emph{dīpa agni}. The toy
corpus has 13 \vk{surf} types but 9 \vk{lem} types. At scale, \vk{lid}
splits \emph{brahman} into the neuter (\emph{the Absolute}) and the
masculine (\emph{priest}).

\paragraph{Ambiguity and OOV.} With $X_k$ the token of a random occurrence
in view $k$, the conditional entropy $H(X_{\vk{lid}}\mid X_{\vk{lem}})$ is
the homography that lemma IDs resolve, $H(X_{\vk{lem}}\mid X_{\vk{seg}})$
the ambiguity a lemmatiser resolves, and $H(X_{\vk{seg}}\mid X_{\vk{lem}})$
the morphology it discards. For OOV, let $\OOV_k$ be the share of benchmark
words with fewer than $\tau$ occurrences in view $k$. Then
$\OOV^{\mathrm{raw}}_{\vk{surf}}=\Delta_{\mathrm{enc}}+\Delta_{\mathrm{sandhi}}
+\Delta_{\mathrm{infl}}+\Delta_{\mathrm{homo}}+\OOV_{\vk{lid}}$, with
$\Delta_{\mathrm{sandhi}}=\OOV_{\vk{surf}}-\OOV_{\vk{seg}}$ and so on: a
telescoping identity that attributes the OOV rate to encoding, sandhi,
inflection and homograph splitting, leaving a residual.

\paragraph{Models as PMI estimators.} All four models estimate (variants of)
the pointwise mutual information matrix,
$\PMI(w,c)=\log\frac{P(w,c)}{P(w)P(c)}$:
\begin{itemize}
\item \textbf{PPMI-SVD}: truncated SVD of the positive, context-smoothed
  PMI matrix \citep{levy2015improving}, the best rank-$d$ approximation
  \citep{eckart1936approximation};
\item \textbf{word2vec} (skip-gram with negative sampling,
  \citealp{mikolov2013distributed}): at its optimum
  $\mathbf{w}^\top\mathbf{c}=\PMI(w,c)-\log k$ \citep{levy2014neural};
\item \textbf{GloVe} \citep{pennington2014glove}: weighted least squares on
  $\log$ co-occurrence counts, a weighted PMI factorisation;
\item \textbf{FastText} \citep{bojanowski2017enriching}: word2vec over
  character $n$-grams, which yields vectors for unseen strings.
\end{itemize}
A word's average PMI with its own contexts is the KL divergence of its
context distribution from the background, so a particle that occurs
everywhere gets small PMI values however frequent it is.

\ex Counting co-occurrence within a sentence, \emph{jala} and the particle
\emph{ca} each co-occur once with each other and \emph{jala} once with
\emph{pā}, yet $\PMI(\textit{jala},\textit{pā})=1.81$ bits and
$\PMI(\textit{jala},\textit{ca})=0.81$: \emph{ca} keeps company with
everything. On positive-PMI rows, cos(\emph{jala}, \emph{toya}) = 0.92 and
cos(\emph{jala}, \emph{agni}) = 0.08 (raw counts: 0.75 and 0.25).

\paragraph{Lookup.} In \vk{lid} a string may name several vectors; we use
MaxSim \citep{reisinger2010multi}, the cosine of the closest pair of
readings.

%==========================================================================
\section{Data and benchmark}\label{sec:data}
%==========================================================================
\subsection{Corpus}\label{sec:corpus}
We parse the DCS CoNLL-U release \citep{hellwig2010dcs} into the five
views, in SLP1 (one ASCII letter per phoneme).

\paragraph{A training-corpus leak.} DCS contains the Amarakośa itself and
eight further synonym lexicons (13,095 sentences). EvalSan's relatedness and
synonym data are derived from the Amarakośa, so a model trained on all of
DCS, EvalSan's included, has read the test dictionary's synonym lists. We
train on DCS without these lexicons (5.71M words); \S\ref{sec:decisions}
shows the size of the effect.

\paragraph{Lemma normalisation.} DCS lemmatises participles to their root
(\emph{gatāḥ} $\to$ \emph{gam}). Pāṇini instead counts \emph{kṛt}-derivatives
and compounds as nominal stems (A.~1.2.46 \emph{kṛt-taddhita-samāsāś ca}).
Following \emph{remove inflection, keep derivation}, profile \vk{np}
applies two rules to every corpus token: anusvāra before a stop becomes the
homorganic nasal (\vk{n}), and participles and gerundives become their
nominal stem (\vk{p}: \emph{gatāḥ} $\to$ \emph{gata}, \emph{pibantaḥ} $\to$
\emph{pibat}), keeping the root's ID in \vk{lid}. The stem rule agrees
96.5\% with the stems DCS itself gives inside compounds. An optional third
rule fuses strongly associated compound members (\vk{c}).

\paragraph{Respelling at lookup.} Some benchmark words follow other citation
conventions than DCS. When a benchmark word is OOV, we try its variants
under optional rules of the grammar: final voicing in pausa (8.2.39,
8.4.56; \emph{kadācit}/\emph{kadācid}), visarga (8.3.15;
\emph{mithaḥ}/\emph{mithas}), anusvāra (8.3.23; \emph{om}/\emph{oṃ}),
nominative vs.\ stem citation (\emph{bhogī}/\emph{bhogin}), and optional
gemination (8.4.46--47; \emph{patra}/\emph{pattra}). Gemination is undone
only where it is optional; intervocalic doubling is lexical
(\emph{kuṭṭaka} $\neq$ \emph{kuṭaka}). This maps 36 benchmark words.

\subsection{Benchmark audit}\label{sec:audit}
\paragraph{Encoding.} SLP1 and WX assign letters differently: \texttt{t} is
dental \emph{t} in SLP1 but retroflex \emph{ṭ} in WX. EvalSan's
relatedness and synonym files are WX, while its categorisation files and
vocabularies are SLP1, and nothing converts between them; the neighbours in
its saved notebook outputs confirm SLP1 vocabularies. Reading each word both
ways and scoring the readings by dictionary lookup and a character-trigram
model shows that no file mixes schemes, so whole-file conversion is the fix.
Only 34\% of relatedness test words are spelled alike in both schemes.

\paragraph{Protocol and defects.} EvalSan gives OOV words the mean vector
(so two OOV words count as synonyms), uses fixed cosine thresholds, and
reports the best categorisation purity of six configurations chosen on test.
Independently of encoding, the relatedness data has 25 self-pairs, 41
duplicates, 14 conflicting labels and 35 dev pairs repeated in test; 112
synonym items repeat a word or have two correct answers. Gold answers are
sound (99.7\% share an Amarakośa meaning with the question). We report on a
cleaned copy (4,509 test and 4,478 dev pairs, 2,922 synonym items).

\subsection{Evaluation protocol}\label{sec:protocol}
Following \citet{schnabel2015evaluation} and \citet{faruqui2016problems}:
OOV items are excluded and \emph{coverage} is reported, and models are
compared on their \emph{common support}. Relatedness is scored by Spearman
$\rho$ and by $\mathrm{AUC}_{ij}=\Pr(s(X_i)>s(X_j))$ for same-synset (2),
same-varga (1) and unrelated (0) pairs, which are invariant to each model's
cosine scale; 3-way thresholds are tuned on dev. Clusterings are scored by
NMI against a permutation null. Everything that is chosen is chosen on
\emph{dev} (relatedness dev split, even-numbered synonym items); test is
used once.

\ex Both toy synonym pairs score 0.92 and the unrelated pairs 0.08 and 0.20,
so AUC$_{20}=1$. But \emph{jala} and \emph{nadī} (``river''), same varga,
share no context and score 0, so AUC$_{10}=0$. Context vectors measure
\emph{paradigmatic} similarity (words that replace each other), while a
varga also groups \emph{syntagmatic} associates (words that occur
together). \S\ref{sec:limits} shows the same effect at scale.

%==========================================================================
\section{Decisions and results}\label{sec:results}
%==========================================================================
In progress: decision log, test results, comparison with Sandhan et al.

\bibliography{refs}

\end{document}
